\documentclass[11pt,a4paper]{article}
\usepackage[margin=24mm]{geometry}
\usepackage{amsmath,amssymb,amsthm,booktabs,longtable,microtype}
\usepackage[T1]{fontenc}
\usepackage{lmodern}
\usepackage[colorlinks=true,linkcolor=blue,urlcolor=blue,citecolor=blue]{hyperref}
\newcommand{\FS}{\operatorname{FS}}
\newcommand{\Z}{\mathbb Z}
\newcommand{\T}{\mathcal T}
\newtheorem{theorem}{Theorem}[section]
\newtheorem{lemma}[theorem]{Lemma}
\newtheorem{proposition}[theorem]{Proposition}
\title{Subset Sums in a Fourfold Sumset\\\large A checked continuation from 103 to 100}
\author{Sum-of-Subset research project}
\date{16 October 2026}
\begin{document}
\maketitle
\begin{abstract}
We continue the repository's independently checked upper bound 103 to 100.
The obstruction at size 103 is a possible empty covering unit. We handle
this obstruction and the corresponding cases at sizes 102 and 101 by
2,580 explicitly specified linear systems. Freshly generated integer
certificates prove every system inconsistent. The supplied verifier uses
only the Python standard library and checks all coefficients exactly.
A uniform effective-hole test treats the high-hole cases, leaving only
one empty unit and, at size 101, at most one outside center.
Together with the known construction this gives the checked interval
$71\le N\le100$. The archived claim $N\le84$ is stronger and is not
certified in its entirety by this work.
\end{abstract}
\tableofcontents

\section{Conclusion, conventions, and dependency boundary}
For a finite real set $A$, let
\[
 \FS(A)=\left\{\sum_{x\in X}x:X\subseteq A\right\},\qquad
 4A=A+A+A+A.
\]
The empty sum belongs to $\FS(A)$, and repetitions are permitted in $4A$.
Write $T_r=r(r+1)/2$ and let $N$ be the maximum feasible cardinality.
\begin{theorem}\label{thm:main}
If $\FS(A)\subseteq4A$, then $|A|\le100$.
\end{theorem}
This report uses the audited 115 foundation and the structural lemmas
proved in the 115-to-103 continuation \cite{previous}. In particular it
uses that report's corrected peeling argument, positive-side stability
and inverse lemmas, common-lattice method, four-extreme estimate, and
active-center identities. It rechecks the size-dependent applicability
gates below; it does not merely assume that a statement proved at 104
also holds at 101. The previous theorem $N\le103$ closes the upper end
of the new descent.

The one-empty-unit strategy and its necessary inequalities already appear
in the archived upper-86 report \cite{archive}. Here we give a compact
fresh verification of the required portion, rebuild the complete model
universe, regenerate all integer certificates, and replace separate
high-hole bootstrap chains by one parameter test. No archived endpoint
certificate or historical PASS record is an input to the new verifier.
This improves the independently closed bound, not the project's strongest
archived numerical claim. It does not determine the exact maximum.

The 71-element example remains
\[
 [-31,31]_{\Z}\cup\{\pm93,\pm155,\pm217,\pm713\}.
\]
Its proof and exact certificate are unchanged. Thus the checked range is
$71\le N\le100$.

\section{Reduction to common-lattice scalar parameters}
Adjoining zero preserves feasibility: $4A\subseteq4(A\cup\{0\})$
and $\FS(A\cup\{0\})=\FS(A)$.
For integralization, fix a four-term representation of every subset sum.
The resulting homogeneous rational linear system has rational solutions
arbitrarily close to its real solution inside its kernel. Preserve the
strict coordinate order and the zero coordinate, then clear denominators.
Consequently existence can be studied with zero present and with integer
coordinates, at each stage of a descending exclusion.

Remove the four most positive and four most negative elements to obtain
$D$, and make two four-extreme peels $D\supset D_1\supset E$ preserving
zero. For $n\in\{101,102,103\}$ put
\[
 d=n-8,\quad e=n-16,\quad E=P\cup\{0\}\cup(-Q),\quad
 a=|P|,\quad b=|Q|,\quad m=a+b=n-17.
\]
The elementary cover and peel comparisons from \cite{previous} give
\[
 |\FS(E)|\le24d+77.
\]
If fewer than eight outer elements were removed, the core is one-sided
with at least $n-16$ nonzero elements. This contradicts
$1+T_{n-16}>24n+77$, checked at all three sizes.
If the core has $p$ opposite pairs and $q$ unpaired absolute values, the
same argument as in \cite{previous} gives
\[
 |\FS(E)|\ge\frac{e^2+3+q^2}{4},\quad
 q^2\le96e+1073-e^2,\quad q\equiv e-1\pmod2.
\]
The resulting lower bounds on $p$ are 22, 21, 20 at sizes 103, 102, 101.
In particular both signs and the required opposite pair are present.

For a positive set $X$ of size $r$, write
\[
 \mu(X)=|\FS(X)|-1-T_r,\quad
 \delta_0(X)=|\{x\in X:2x\notin R(X\cup\{0\})\}|,\quad
 \rho(X)=\mu(X)+2-2\delta_0(X),
\]
where $R(X)$ denotes restricted sums of two distinct elements. The
inherited estimates are $\rho(X)\ge0$ for $r\ge14$, and
$\mu(X)\ge r-3$, $\rho(X)\ge r-14$ when its support is not a progression.

Set $\Lambda=1+T_a+T_b$ and $\Omega=22d+97-\Lambda$.
The cover slack is
\[
 c=2|2D|+18d+85-|\FS(D)|\ge0.
\]
With $s=|\FS(E)|$ and $\mu=\mu(P)+\mu(Q)$, let
$\kappa=s-\Lambda-\mu$. If $\ell_1,\ell_2$ are the two nonnegative
peeling slacks and
$g=4d-10-(|R(D)|-|R(D_1)|)$, set
$\T=\kappa+\ell_1+\ell_2+g$. The inherited exact identity is
\begin{equation}\label{eq:budget}
 \rho(P)+\rho(Q)+2\theta+c+\T=\Omega,
 \qquad \theta=\delta_0(P)+\delta_0(Q)+8-\Delta(D)\ge0.
\end{equation}
Here $\Delta(D)=|2D|-|R(D)|$.

The two-set inverse theorem and propagation to $D$ require precisely the
strict gates
\begin{equation}\label{eq:gates}
 T_{\min(a,b)}-2-\Omega>0,\quad
 2m+4e+14-\Lambda+3\Omega<0,\quad
 \Lambda-(4d-10+\Omega)>0.
\end{equation}
Enumerate both orientations of all $a,b$ at least the paired minimum,
with $a+b=m$ and $\Omega\ge0$. The margins are:
\begin{center}
\begin{tabular}{rrrrrr}
\toprule
$n$&Orientations&Paired minimum&First minimum&Second maximum&Third minimum\\
\midrule
103&35&22&335&$-477$&1229\\
102&36&21&299&$-377$&1169\\
101&37&20&264&$-277$&1109\\
\bottomrule
\end{tabular}
\end{center}
All hypotheses in the common-lattice argument of \cite{previous} are
therefore valid. After scaling, $D$ and $E$ are integral. If $h_P,h_Q,z,w$
are the hull-hole counts of $\FS(P),\FS(Q),\FS(E),\FS(D)$ on this lattice,
then
\begin{gather}
 h_P,h_Q\le\kappa,\quad z=h_P+h_Q-\kappa,
 \quad0\le z\le\min(h_P,h_Q)\le\kappa,\quad
 h_X<|\FS(X)|-1,\label{eq:holes}\\
 w\le4z,\qquad
 \T\ge\max\{z,G-2z,t_0\},\quad
 G=2d-12,\quad t_0=\lceil G/3\rceil.\label{eq:T}
\end{gather}
For $z>0$ neither side support is a progression, so
$\rho(P)+\rho(Q)\ge m-28$. All relevant side sizes are at least 20,
above the inherited defect and span thresholds.

\section{One effective-hole test for all high-hole states}
The counted endpoint lemma of \cite{previous} says that if a positive
prefix of size $i$ has excess increment $y_i=t\ge0$ and $i\ge2t+5$,
its support, after scaling, is
\[
 [0,L]_{\Z}\setminus(H\cup(L-H)),\qquad
 H\subseteq[1,2t]_{\Z},\quad |H|\le t.
\]
It follows from the small-increment inverse lemma and Lev's dense
subset-sum interval theorem as proved there. The following parameterized
version makes the finite high-hole audit particularly short.

\begin{lemma}\label{lem:effective}
Let $C\ge0$, $\beta\ge1$ be integers and define
\begin{gather*}
 K=\max\{k\ge3:k(k-3)/2\le C+1\},\quad D_*=C+K,\quad E_*=D_*+1,\\
 M=\max\left\{2\beta,1+\left\lfloor
 \frac{1+\sqrt{9+8(E_*+6\beta)}}2\right\rfloor\right\},
 \quad J=\max\{K+1,2\beta+5,M\}.
\end{gather*}
Suppose
\begin{equation}\label{eq:effective-tests}
 \beta(r-J+1)>D_*,\qquad T_{J-1}-1>D_*.
\end{equation}
If a positive $r$-element set has, on a specified lattice,
$2\beta<h<|\FS|-1$ hull holes, then $\mu\ge r+C+1$.
\end{lemma}
\begin{proof}
Assume $\mu\le r+C$ and divide by the specified lattice step. The set
is integral and primitive, since a gcd at least two would create at
least $|\FS|-1$ holes. Its support is not a progression.
The last progression prefix of size $k\ge3$, if there is one, is
followed by a disjoint translate. Thus
$\mu\ge r-1+k(k-3)/2$ and $k\le K$. If there is no such prefix, use
the tail starting at size four. In either case all increments in a tail
starting no later than $K+1$ obey
\[
 y_j\ge1,\qquad \sum(y_j-1)\le D_*.
\]
The first test forces an $i\in[J,r]$ with $1\le y_i\le\beta$.
The counted endpoint lemma gives a scaled support with
$H\subseteq[1,2\beta]$ and $|H|\le\beta$. Its primitive coefficient
set makes its scale an integer. If the scale exceeded one, the first
later off-scale element would give
\[
 y_j-1\ge T_{j-1}-j\ge T_{i-1}-1\ge T_{J-1}-1>D_*,
\]
a contradiction. The scale is therefore one.

Every later increment excess is at most $E_*$. At a prefix of size
$v\ge M$, its total $L\ge T_v$ and its next element $x>v\ge2\beta$.
If $x\ge L-4\beta$, the new translate contributes at least $L-6\beta$
points above $L$. This exceeds $v+1+E_*$, since the definition of $M$
gives $T_M-6\beta>M+1+E_*$ and the difference increases with $v$.
Hence $2\beta<x<L-4\beta$. The two central intervals overlap or abut;
the lower holes remain $H$, and symmetry gives the upper holes.
Induction leaves at most $2\beta$ holes, the required contradiction.
\end{proof}

Here is the complete finite rule used by the verifier. For
$20\le r\le64$, $7\le\beta\le20$, $0\le C\le600$, retain the rows
satisfying \eqref{eq:effective-tests}. The largest valid $C$ for a pair
$(r,\beta)$ is only a convenient witness, not an asserted global optimum.
For each $z\ge15$, use only rows with $2\beta<z$. Recheck the chosen
row exactly, and put
\[
 u_r=\max\bigl(\{r-3\}\cup\{r+C+1:\text{eligible rows}\}\bigr),
 \qquad v_r=\max\{r-14,u_r+2-2r\}.
\]
Equation \eqref{eq:holes} licenses the test for both sides. The second
bound uses $\delta_0\le r$. Equation \eqref{eq:budget} implies
\begin{equation}\label{eq:cmax}
 c\le c_{\max}:=\Omega-v_a-v_b-\max\{z,G-2z,t_0\}.
\end{equation}
Also $0\le z\le\kappa\le\T\le\Omega$, so this is a finite enumeration
with a proved upper limit. A negative $c_{\max}$ discards a state.

\section{Coverage inequalities and all-active cases}
Write $R_p=\max P$, $R_q=\max Q$, $M_p=\max D$, $M_q=-\min D$.
Let $S_p,S_q$ be the sums of the positive and negative magnitudes of $D$,
and put
\[
 \varepsilon_p=\sum P-T_a=\mu(P)+h_P,\quad
 \varepsilon_q=\sum Q-T_b=\mu(Q)+h_Q,
\]
\[
 F=S_p+S_q+c-2(M_p+M_q)-2(R_p+R_q).
\]
Only the ten positive-positive and ten negative-negative pairs of the
four peeled elements of each sign can lie outside
$[-M_q-R_q,M_p+R_p]$. Thus
\begin{equation}\label{eq:Fupper}
 |2D|\le M_p+M_q+R_p+R_q+21,\qquad
 F\le18d+126+w\le18d+126+4z.
\end{equation}
The outer magnitudes are denoted $p_1>p_2>p_3>p_4>M_p$ and
$q_1>q_2>q_3>q_4>M_q$. Define
\[
 \lambda_p=p_1-p_2-p_3-p_4,\quad
 \alpha_p=p_2-p_3-p_4,\quad
 \beta_p=2p_2-p_1-p_3-p_4,
\]
and their reflected counterparts. The ten ordinary centers on each side
are indexed as
\begin{align*}
 v^p&=(\lambda_p+p_1,\ldots,\lambda_p+p_4,
          \lambda_p-q_1,\ldots,\lambda_p-q_4,\alpha_p,\beta_p),\\
 v^q&=(\lambda_q+q_1,\ldots,\lambda_q+q_4,
          \lambda_q-p_1,\ldots,\lambda_q-p_4,\alpha_q,\beta_q).
\end{align*}
Indices are $0,\ldots,9$. The first paired unit consists of $p4,q4$.
In branch A, $\beta_p+\beta_q\le0$, the second is $p9,q9$.
In branch B, $\beta_p+\beta_q>0$, the second is $p5,q5$.
The other sixteen units are single centers. Each of the eighteen units
has capacity $d$ in the inherited cover; an empty unit costs $d$ in $c$.
Losses are counted by formal slots even if numerical centers coincide.

If all units are nonempty, the active-center proof in \cite{previous}
applies, regardless of whether $c<d$. After a possible reflection it gives
one of
\begin{align}
 F&\ge F_A(a,b)+\tfrac12\varepsilon_p+\varepsilon_q,
 &F_A&=\tfrac12T_a+T_b+6a+8b+85,\label{eq:FA}\\
 F&\ge F_B(a,b)+\tfrac23\varepsilon_p+\varepsilon_q,
 &F_B&=\tfrac23T_a+T_b+6a+\tfrac{22b+236}{3}.\label{eq:FB}
\end{align}
The proof absorbs all out-of-hull corrections into the available spans.
Its hypothesis is actual activity; $c<d$ is a sufficient condition for it.

For $0\le z\le14$, use $u_a=u_b=0$, $v_a=v_b=0$ at $z=0$;
otherwise use $u_a=a-3$, $u_b=b-3$, $v_a=a-14$, $v_b=b-14$.
For $z\ge15$, use Section 3. In both cases
$\varepsilon_p\ge u_a+z$, $\varepsilon_q\ge u_b+z$ (at zero holes
this is simply nonnegativity). Subtracting \eqref{eq:Fupper} from both
\eqref{eq:FA} and \eqref{eq:FB} gives the checked margins:
\begin{center}\small
\begin{tabular}{rrrrrrr}
\toprule
$n$&Low states&Low margin&High states&High margin&High $c_{\max}$&Empty-eligible low states\\
\midrule
103&261&$215/2$&0&--&--&9\\
102&300&$155/2$&32&365&6&66\\
101&333&$113/2$&399&248&36&165\\
\bottomrule
\end{tabular}
\end{center}
States are necessary scalar relaxations, not constructed sets. All margins
are strictly positive. Every surviving high-hole state has $c_{\max}<d$,
so all high-hole cases are excluded. Every low-hole case with no empty
unit is excluded as well. Only the empty-eligible low states remain.

\section{Exactly one empty unit and at most one outside center}
For low holes the basic bound is
\[
 c\le\Omega-G\quad(z=0),\qquad
 c\le\Omega-G+2z-(m-28)\quad(1\le z\le14).
\]
At the three sizes it is always less than $2d$, hence an empty-unit case
has exactly one such unit. Put $\delta=c-d\ge0$.

\subsection{Outer integrality and exhaustive contributor choices}
Every contributing single center is integral, since it translates the
integral set $D$ into the integral target support. If all four own centers
of a side remain, the identity
$\lambda_p=-v^p_0+v^p_1+v^p_2+v^p_3$ gives integrality of that side.
If one own center on the positive side is missing, all negative own
centers remain. Thus all $q_i$ are integral. The always unpaired positive
centers $\lambda_p-q_3,\lambda_p-q_4$ remain and force $\lambda_p$ integral.
The other three $p_i$ and the defining identity force the last $p_i$
integral. Reflection covers the other case. All outer strict gaps are
therefore at least one.

The first pair always has at most one contributing member: its center
sum is less than $-(M_p+M_q)$. If present, reflect to retain $p4$.
In branch B the second pair also has at most one contributor. In branch A
this holds if both unpaired centers $p5,q5$ remain, because
\[
 \beta_p+\beta_q=-(v^p_5+v^q_5)-2(p_3+p_4+q_3+q_4)
 <-(M_p+M_q).
\]
If $p5$ or $q5$ is the missing unit in branch A, omit the entire second
pair from the retained list. Allowing it to split only weakens the model.

The complete enumeration is now explicit. Choose a branch and a missing
unit among the two pairs and sixteen singles. Retain all other singles.
Retain $p4$ unless the first pair is missing. For the second pair retain
each of its two possible contributors in a separate mode, except when
that pair is missing or deliberately omitted. There are 33 modes in A
and 35 in B. Sixty-six modes retain 17 centers; the two split-pair modes
retain 16. If the first pair is missing both choices for the second are
still included. This loses no configuration.

\subsection{Hull losses and the five regimes}
For a retained positive center $v$, let $l_v,u_v$ be the counts of $v+D$
below 1 and above $S_p$. Reflect for a negative center. Then
\begin{equation}\label{eq:loss}
 \sum_{v\text{ retained}}(l_v+u_v)\le\delta.
\end{equation}
Within its hull a positive center satisfies
\begin{equation}\label{eq:normal}
 -v-l_v\le-(b+5),\qquad v-u_v-S_p\le-(a+4).
\end{equation}
Outside, the valid replacements are
\begin{align}
 v<0:&\quad -v-l_v-M_p\le-d,\quad l_v\ge b+5,\label{eq:outside-low}\\
 v>S_p:&\quad v-u_v-S_p-M_q\le-(d-1),\quad u_v\ge a+5.\label{eq:outside-high}
\end{align}
These follow by counting integer points in the occupied positive and
negative spans; equivalently substitute
$U=M_p-a-4$, $V=M_q-b-4$ in the global quantile bounds of
\cite{previous}. The opposite endpoint inequality in \eqref{eq:normal}
remains valid. Thus each outside center costs at least $\min(a,b)+5$.

For 103, write $a,b=43\mp k$ up to reflection. An empty unit gives
$\delta\le21-k^2$, $k\le4$, and $z=0$; all retained centers are in-hull.
For 102, write $a,b=42-k,43+k$. At zero holes
$\delta\le45-k(k+1)$, $k\le6$, and this is less than $47-k$.
At positive low holes $\delta\le16-k(k+1)$, an even smaller bound.
All retained centers are again in-hull.

For 101, write $a,b=42\mp k$. At positive low holes,
$\delta\le41-k^2<47-k$ and $k\le6$. The core masses satisfy
$\varepsilon_p\ge a-2$, $\varepsilon_q\ge b-2$.
At zero holes, $\delta\le69-k^2$, $k\le8$. Since
\[
 2(47-k)-(69-k^2)=(k-1)^2+24>0,
\]
at most one retained center is outside its hull. If one is outside,
$k^2-k\le22$, hence $k\le5$ and $37\le a,b\le47$.

All remaining configurations are covered by these relaxed regimes,
with $a=a_0+t$, $b=b_0-t$, $0\le t\le t_{\max}$:
\begin{center}\small
\begin{tabular}{lrrrrrrrl}
\toprule
Regime&$n$&$a_0$&$b_0$&$t_{\max}$&$\delta_{\max}$&Upper $F$&Modes&Retained centers\\
\midrule
N103&103&39&47&8&21&1836&68&All in-hull\\
N102&102&36&49&13&45&1874&68&All in-hull\\
Z&101&34&50&16&69&1800&68&All in-hull\\
P&101&36&48&12&41&1856&68&All in-hull; mass bonus\\
E&101&37&47&10&69&1800&2308&One designated exception\\
\bottomrule
\end{tabular}
\end{center}
Regime E enumerates both endpoints of every retained center in every
mode, giving $2(66\cdot17+2\cdot16)=2308$ systems. No guessed exception
location is used. The scalar verifier also checks every state against
these regime ranges and strict outside-loss thresholds.

\section{Complete linear systems and exact certificates}
All variables
\[
 x=(p_1,\ldots,p_4,q_1,\ldots,q_4,M_p,M_q,R_p,R_q,S_p,S_q,
       \delta,t,(l_v,u_v)_{v\text{ retained}})
\]
are nonnegative real numbers. We deliberately relax $t$ and all losses
to real variables. Use the following rows, all oriented $Ax\le b$:
\begin{center}\small
\begin{tabular}{ll}
\toprule
Row&Constraint (and reflected version)\\
\midrule
$D_i^p$ ($i=1,2,3$)&$p_{i+1}-p_i\le-1$\\
$D_4^p$&$M_p-p_4\le-1$\\
$E_p$&$R_p-M_p\le-4$\\
$G_p,G_q$&$-R_p+t\le-a_0$, $-R_q-t\le-b_0$\\
$H_p$&$2M_p+6R_p-2S_p+s_0t\le-(s_0a_0-j_0+12)$\\
$H_q$&$2M_q+6R_q-2S_q-s_0t\le-(s_0b_0-j_0+12)$\\
cap, range&$\delta\le\delta_{\max}$, $t\le t_{\max}$\\
branch&$\beta_p+\beta_q\le0$ (A), $-\beta_p-\beta_q\le-1$ (B)\\
$L_v,U_v$&Equation \eqref{eq:normal}, substituting $a=a_0+t$, $b=b_0-t$\\
loss&Equation \eqref{eq:loss}\\
Fupper&$S_p+S_q+\delta-2(M_p+M_q+R_p+R_q)\le F_{\max}-d$\\
\bottomrule
\end{tabular}
\end{center}
For N103 use $(s_0,j_0)=(87,1849)$; for N102, Z, E use $(85,1764)$;
for P use $(87,1768)$. The mass rows follow from
\[
 S_p\ge T_a+\varepsilon_p+3R_p+6+M_p
\]
and its reflection, together with the exact tangent inequalities
\[
 2T_r-(87r-1849)=(r-43)^2\ge0,\qquad
 2T_r-(85r-1764)=(r-42)^2\ge0.
\]
In P add $2\varepsilon_p\ge2a-4$ and its reflection, giving
$2(T_r+r-2)\ge87r-1768$.
In E replace only the designated $L_v$ or $U_v$ by
\eqref{eq:outside-low} or \eqref{eq:outside-high}, and add its unavoidable
loss row. The opposite endpoint row is unchanged. It is unnecessary to
impose the actual outside location, because omitting it enlarges the model.

\begin{lemma}[Integer contradiction rule]\label{lem:certificate}
If $x\ge0$, $Ax\le b$ and a nonnegative integer row vector $w$ satisfies
\[
 wA\ge0\text{ coefficientwise},\qquad wb<0,
\]
the system is inconsistent.
\end{lemma}
\begin{proof}
One would have $0\le wAx\le wb<0$.
\end{proof}
The file \path{data/upper100_certificates.json} contains 189 distinct
integer row-weight dictionaries and an assignment for every one of the
2,580 model keys. Row labels are exactly those in
\path{verification/upper100.py}; for example, \texttt{Lp5} denotes the
lower row for the positive center with index 5, and \texttt{outside}
denotes the extra unavoidable-loss row in E.
The verifier reconstructs all models from the rules above, requires exact
equality between the generated key set and the assigned key set, checks
that every referenced row exists, and checks Lemma~\ref{lem:certificate}
using integers. Every system passes. The minimum scale-independent
contradiction gap $-wb/\sum_iw_i$ is $8/101>0$.

These are finite proof data, not optimizer verdicts. The optional
discovery script uses SciPy to propose weights, converts them to rational
numbers and then integers, and refuses to save any vector that fails the
exact rule. Neither SciPy nor floating-point arithmetic is used by the
verification script. The freshly generated vectors need not agree with
the archived multipliers.

\section{Closing the descent and reproducing the result}
\begin{proof}[Proof of Theorem~\ref{thm:main}]
Start with $N\le103$ from \cite{previous}. A 103-element feasible set
must contain zero, since otherwise adjoining zero contradicts that bound.
Integralize and use the checked lattice gates. Sections 3--4 exclude
all high-hole states and all low-hole states with no empty unit.
The remaining states have one empty unit and belong to N103, whose
68 exact linear contradictions exclude them. Thus $N\le102$.

At size 102 the same argument applies, and the N102 systems exclude the
remaining empty-unit states. Thus $N\le101$. At size 101, the remaining
states are partitioned into positive low holes (P), zero holes with all
retained centers in-hull (Z), and zero holes with one retained exception
(E). All their systems are inconsistent. Therefore $N\le100$.
At no stage is feasibility assumed to pass to arbitrary subsets.
\end{proof}

From the repository root, the verification commands are
\begin{verbatim}
python -I -S -B verification/upper103.py --output .work/base103.json
python -I -S -B verification/upper100.py --output .work/upper100.json
\end{verbatim}
Python 3.10 or later suffices; no third-party package or network is needed.
The new complete result and receipt are in
\path{results/2026-10-16/upper100/}. A second run with \texttt{-O} produces
the identical JSON. Checks use explicit exceptions, not Python assertions.

Seven negative controls reject an invalid hole witness, zero weights,
a noncontradictory positive certificate, a negative weight, an absent row,
a missing case, and an attempt to certify cardinality 100 using this
three-endpoint routine. Hashes in the output bind the run to the verifier
and certificate bytes. These checks do not replace the written proof that
actual configurations satisfy the rows, nor do they formally certify the
external structural theorems inherited from \cite{previous}.

Optional regeneration, requiring SciPy, is
\begin{verbatim}
python experiments/generate_upper100_certificates.py \
  --output .work/regenerated_upper100.json
python verification/upper100.py \
  --certificates .work/regenerated_upper100.json
\end{verbatim}
The next endpoint has a concrete obstruction to the current one-empty
argument. At $n=100$, $a=41$, $b=42$, $z=0$, one has
\[
 d=92,\quad\Omega=356,\quad G=172,\quad c\le184=2d.
\]
Two empty units are therefore allowed by this relaxation. Excluding them,
or strengthening the budget to remove them, is a necessary next step.
This is an obstruction in the proof method, not a feasible 100-element set.

\begin{thebibliography}{9}
\bibitem{previous} Sum-of-Subset project,
\emph{Subset Sums in a Fourfold Sumset: A checked continuation from the
upper bound 115 to 103}, 14 October 2026.
\path{reports/upper103_from115_20261014.tex}.
Its Sections 1--6 specify the audited foundation, external theorems,
and structural lemmas used here.
\bibitem{archive} Sum-of-Subset project,
\emph{Complete upper-86 report}, 5 October 2026.
\path{reports/upper86_complete_20261005.tex}, sections on upper bounds
102 and 100 and their certificate appendices. Source of the empty-unit
strategy; its complete endpoint chain is outside this audit.
\end{thebibliography}
\end{document}
